\documentclass[11pt,a4paper]{article}

\usepackage[T1]{fontenc}
\usepackage[utf8]{inputenc}
\usepackage{lmodern}
\usepackage{geometry}
\usepackage{hyperref}
\usepackage{xcolor}
\usepackage{listings}
\usepackage{enumitem}
\usepackage{longtable}
\usepackage{setspace}

\geometry{margin=1in}
\setstretch{1.12}

\hypersetup{
  colorlinks=true,
  linkcolor=blue!60!black,
  urlcolor=blue!60!black,
  pdftitle={Distributed Sales Synchronization - Implementation Report},
  pdfauthor={Project Team}
}

\definecolor{codebg}{RGB}{248,248,248}
\definecolor{codekw}{RGB}{0,70,140}
\definecolor{codecom}{RGB}{0,120,0}
\definecolor{codestr}{RGB}{170,60,0}

\lstdefinestyle{javastyle}{
  language=Java,
  backgroundcolor=\color{codebg},
  basicstyle=\ttfamily\small,
  keywordstyle=\color{codekw}\bfseries,
  commentstyle=\color{codecom},
  stringstyle=\color{codestr},
  showstringspaces=false,
  breaklines=true,
  frame=single,
  framerule=0.2pt,
  tabsize=2
}

\lstdefinestyle{propstyle}{
  language={},
  backgroundcolor=\color{codebg},
  basicstyle=\ttfamily\small,
  showstringspaces=false,
  breaklines=true,
  frame=single,
  framerule=0.2pt,
  tabsize=2
}

\title{Distributed Sales Synchronization\\Implementation and Design Report}
\author{DB\_synchronisation Project}
\date{\today}

\begin{document}
\maketitle
\tableofcontents
\newpage

\section{Executive Summary}
This project implements a distributed data synchronization platform for two branch-office systems (BO1 and BO2) and a head-office system (HO). The two branch systems are independent operational systems with their own local PostgreSQL databases and local web frontends for CRUD operations. The HO system is a consolidation platform that receives events over RabbitMQ and materializes a central \texttt{consolidated\_sales} dataset.

The central architectural idea is \textbf{event-driven synchronization}: every branch-side create, update, and delete operation emits a RabbitMQ message, and HO consumes these messages to perform idempotent upsert/delete operations in its own database.

\section{Problem Context and Goals}
\subsection{Business Context}
The system models geographically separated sales branches that need to continuously synchronize operational data to a central office. The architecture must satisfy the following:
\begin{itemize}[leftmargin=*]
  \item Allow branches to operate locally and independently.
  \item Keep HO in sync with branch data.
  \item Support insert, update, and delete semantics end-to-end.
  \item Avoid strong temporal coupling between branch request execution and HO data materialization.
\end{itemize}

\subsection{Technical Goals}
\begin{itemize}[leftmargin=*]
  \item Expose standalone BO and HO web applications with frontend pages.
  \item Expose CRUD REST APIs on all services.
  \item Use RabbitMQ as asynchronous transport from BO to HO.
  \item Use durable messaging and resilient exchange/queue configuration.
  \item Preserve clear message contract and deterministic HO merge logic.
\end{itemize}

\section{System Architecture}
\subsection{High-Level Flow}
\begin{enumerate}[leftmargin=*]
  \item A user creates/updates/deletes a sale in BO1 or BO2 UI.
  \item BO persists the local transaction in its \texttt{product\_sales} table.
  \item BO publishes a domain event to RabbitMQ with operation type \texttt{UPSERT} or \texttt{DELETE}.
  \item HO listens to the queue, deserializes the event, and applies the corresponding mutation to \texttt{consolidated\_sales}.
\end{enumerate}

\subsection{Module Roles}
\begin{longtable}{|p{0.2\linewidth}|p{0.72\linewidth}|}
\hline
\textbf{Module} & \textbf{Role} \\
\hline
BO1 / BO2 & Local operational systems: CRUD API + frontend + direct event publisher to RabbitMQ. \\
\hline
HO & Consolidation system: queue listener + idempotent upsert/delete + consolidation API + frontend. \\
\hline
RabbitMQ & Message broker for asynchronous BO $\rightarrow$ HO transport. \\
\hline
PostgreSQL & Persistence layer for each module (\texttt{bo1\_db}, \texttt{bo2\_db}, \texttt{ho\_db}). \\
\hline
\end{longtable}

\section{Data Contract and Event Semantics}
\subsection{Message Structure}
The message schema is represented by \texttt{SaleMessage}. It includes:
\begin{itemize}[leftmargin=*]
  \item \texttt{operation} (\texttt{UPSERT} or \texttt{DELETE})
  \item \texttt{boId} and \texttt{localSaleId} as event identity across modules
  \item Sales payload fields (date, region, product, qty, cost, amt, tax, total)
\end{itemize}

\begin{lstlisting}[style=javastyle,caption={Simplified event contract (BO side)}]
public class SaleMessage {
    private String operation;
    private String boId;
    private Long localSaleId;
    private LocalDate date;
    private String region;
    private String product;
    private Integer qty;
    private BigDecimal cost;
    private BigDecimal amt;
    private BigDecimal tax;
    private BigDecimal total;
}
\end{lstlisting}

\subsection{Why an Explicit Operation Field?}
The \texttt{operation} field avoids ambiguity in event interpretation:
\begin{itemize}[leftmargin=*]
  \item \texttt{UPSERT} means ``create or update this logical sale key''.
  \item \texttt{DELETE} means ``remove this logical sale key''.
\end{itemize}
This design keeps the consumer logic deterministic and explicit, and avoids relying on null field heuristics to infer intent.

\section{RabbitMQ Design and Rationale}
\subsection{Exchange and Queue Topology}
The system uses:
\begin{itemize}[leftmargin=*]
  \item \textbf{Topic exchange}: \texttt{sales.exchange} (durable)
  \item \textbf{Queue}: \texttt{sales.hq} (durable)
  \item \textbf{Routing key}: \texttt{sales.hq}
\end{itemize}

\begin{lstlisting}[style=javastyle,caption={HO queue/exchange/binding configuration}]
@Bean
public Queue salesQueue(@Value("${rabbitmq.queue}") String queueName) {
    return new Queue(queueName, true);
}

@Bean
public TopicExchange salesExchange(@Value("${rabbitmq.exchange}") String exchangeName) {
    return new TopicExchange(exchangeName, true, false);
}

@Bean
public Binding salesBinding(Queue salesQueue, TopicExchange salesExchange,
                            @Value("${rabbitmq.routing-key}") String routingKey) {
    return BindingBuilder.bind(salesQueue).to(salesExchange).with(routingKey);
}
\end{lstlisting}

\subsection{Durability Choices}
Two durability decisions are implemented:
\begin{enumerate}[leftmargin=*]
  \item Durable exchange and durable queue declarations.
  \item Persistent message delivery mode at publish time.
\end{enumerate}

\begin{lstlisting}[style=javastyle,caption={Publisher uses persistent delivery mode}]
public void publishEvent(SaleMessage saleMessage) {
    rabbitTemplate.convertAndSend(exchangeName, routingKey, saleMessage, message -> {
        message.getMessageProperties().setDeliveryMode(MessageDeliveryMode.PERSISTENT);
        return message;
    });
}
\end{lstlisting}

This combination improves survivability in broker restarts and aligns with the reliability requirements of distributed synchronization.

\subsection{JSON Serialization Strategy}
The project configures Jackson-based message conversion and registers \texttt{JavaTimeModule} to correctly serialize/deserialize \texttt{LocalDate}.

\begin{lstlisting}[style=javastyle,caption={Message converter with Java time support}]
@Bean
public MessageConverter jacksonMessageConverter() {
    ObjectMapper objectMapper = new ObjectMapper();
    objectMapper.registerModule(new JavaTimeModule());
    return new Jackson2JsonMessageConverter(objectMapper);
}
\end{lstlisting}

Rationale:
\begin{itemize}[leftmargin=*]
  \item Human-readable payloads for observability/debugging.
  \item Strong compatibility with Spring AMQP and Java domain objects.
  \item Correct support for date/time types.
\end{itemize}

\section{Producer Flow (BO1/BO2)}
\subsection{Current Implementation: Direct Publish per CRUD}
Each BO service now publishes immediately after the local database mutation in the service layer.

\begin{lstlisting}[style=javastyle,caption={BO direct publish on create/update/delete}]
@Transactional
public ProductSale createSale(SaleUpsertRequest request) {
    ProductSale sale = new ProductSale();
    applyRequest(sale, request);
    ProductSale saved = productSaleRepository.save(sale);
    publishUpsert(saved);
    return saved;
}

@Transactional
public void deleteSale(Long id) {
    ProductSale sale = productSaleRepository.findById(id).orElseThrow();
    productSaleRepository.delete(sale);
    publishDelete(id);
}
\end{lstlisting}

\subsection{Why We Switched to Direct Publish}
The earlier design used a scheduler and an outbox table. That flow generated repeated polling queries and periodic logging noise. The redesign objective was to simplify operational behavior and remove unnecessary periodic Hibernate select statements.

Reasons for switching:
\begin{itemize}[leftmargin=*]
  \item Lower implementation complexity for this project stage.
  \item Immediate event emission, reducing perceived sync latency.
  \item Cleaner BO logs (no periodic outbox polling query flood).
\end{itemize}

Trade-off acknowledged:
\begin{itemize}[leftmargin=*]
  \item Direct publish reduces decoupling between DB write and event publish compared to a transactional outbox pattern.
\end{itemize}

\section{Consumer Flow (HO)}
\subsection{Listener Behavior}
HO receives events and applies operation-specific logic.

\begin{lstlisting}[style=javastyle,caption={HO listener applies DELETE vs UPSERT semantics}]
@RabbitListener(queues = "${rabbitmq.queue}")
@Transactional
public void consumeSale(SaleMessage message) {
    if ("DELETE".equalsIgnoreCase(message.getOperation())) {
        consolidatedSaleRepository.deleteByBoIdAndLocalSaleId(
            message.getBoId(), message.getLocalSaleId());
        return;
    }

    ConsolidatedSale consolidatedSale = consolidatedSaleRepository
        .findByBoIdAndLocalSaleId(message.getBoId(), message.getLocalSaleId())
        .orElseGet(ConsolidatedSale::new);

    if (consolidatedSale.getId() == null) {
        consolidatedSale.setBoId(message.getBoId());
        consolidatedSale.setLocalSaleId(message.getLocalSaleId());
    }

    mapFields(consolidatedSale, message);
    consolidatedSaleRepository.save(consolidatedSale);
}
\end{lstlisting}

\subsection{Idempotence and Merge Key}
The logical identity is the pair \texttt{(boId, localSaleId)}. It is enforced both in logic and in schema via a unique constraint.

\begin{lstlisting}[style=javastyle,caption={HO consolidated entity unique key}]
@Entity
@Table(
  name = "consolidated_sales",
  uniqueConstraints = @UniqueConstraint(columnNames = {"bo_id", "local_sale_id"})
)
public class ConsolidatedSale { ... }
\end{lstlisting}

This design prevents duplicate consolidated rows and makes repeated UPSERT events safe.

\section{Configuration and Deployment Notes}
\subsection{Runtime Configuration}
BO and HO are parameterized through \texttt{.env} files imported by Spring.

\begin{lstlisting}[style=propstyle,caption={BO relevant properties}]
spring.config.import=optional:file:.env[.properties]
server.port=${SERVER_PORT}
spring.rabbitmq.host=${RABBITMQ_HOST}
spring.rabbitmq.port=${RABBITMQ_PORT}
rabbitmq.exchange=${RABBITMQ_EXCHANGE}
rabbitmq.routing-key=${RABBITMQ_ROUTING_KEY}
bo.id=${BO_ID}
\end{lstlisting}

\begin{lstlisting}[style=propstyle,caption={HO relevant properties}]
spring.config.import=optional:file:.env[.properties]
server.port=${SERVER_PORT}
spring.rabbitmq.host=${RABBITMQ_HOST}
spring.rabbitmq.port=${RABBITMQ_PORT}
rabbitmq.queue=${RABBITMQ_QUEUE}
rabbitmq.exchange=${RABBITMQ_EXCHANGE}
rabbitmq.routing-key=${RABBITMQ_ROUTING_KEY}
\end{lstlisting}

\subsection{Why Frontend Does Not Need Separate Docker}
Each frontend is served as static content by its Spring Boot service itself:
\begin{itemize}[leftmargin=*]
  \item BO1 UI from BO1 server.
  \item BO2 UI from BO2 server.
  \item HO UI from HO server.
\end{itemize}
So there is no independent frontend container requirement.

\section{Design Decisions: Detailed Justification}
\subsection{Decision 1: Event-Driven over Direct Synchronous BO-to-HO API}
\textbf{Chosen:} RabbitMQ asynchronous events.\\
\textbf{Why:}
\begin{itemize}[leftmargin=*]
  \item Better temporal decoupling: BO does not block waiting for HO API processing.
  \item Better scalability for additional branch systems.
  \item Better resilience characteristics through durable broker constructs.
\end{itemize}

\subsection{Decision 2: Topic Exchange Topology}
\textbf{Chosen:} topic exchange with explicit routing key.\\
\textbf{Why:}
\begin{itemize}[leftmargin=*]
  \item Flexible future extension for more consumers/routing patterns.
  \item Compatible with current simple single-route setup.
\end{itemize}

\subsection{Decision 3: UPSERT/DELETE Event Model}
\textbf{Chosen:} operation-aware messages.\\
\textbf{Why:}
\begin{itemize}[leftmargin=*]
  \item Mirrors CRUD semantics exactly.
  \item Avoids overloading one event shape for all intents.
  \item Simplifies consumer branch logic and testing.
\end{itemize}

\subsection{Decision 4: Composite Business Key in HO}
\textbf{Chosen:} unique key \texttt{(boId, localSaleId)}.\\
\textbf{Why:}
\begin{itemize}[leftmargin=*]
  \item Natural identity in a multi-branch distributed context.
  \item Enables deterministic merge and idempotent reprocessing.
\end{itemize}

\subsection{Decision 5: Direct Publish Redesign}
\textbf{Previous:} scheduler + outbox polling model.\\
\textbf{Current:} direct publish per CRUD call.\\
\textbf{Why current decision:}
\begin{itemize}[leftmargin=*]
  \item Remove continuous query noise from polling.
  \item Reduce conceptual complexity for the current scope.
  \item Provide immediate propagation behavior.
\end{itemize}

\section{Operational Behavior of RabbitMQ in This System}
\subsection{Lifecycle of a Single CRUD Action}
Assume BO1 updates one sale row:
\begin{enumerate}[leftmargin=*]
  \item Request reaches BO1 REST endpoint.
  \item Service updates \texttt{product\_sales} in BO1 PostgreSQL.
  \item Service creates a \texttt{SaleMessage(operation=UPSERT, boId=BO1, localSaleId=...)}.
  \item \texttt{RabbitTemplate} publishes to \texttt{sales.exchange} with routing key \texttt{sales.hq}.
  \item RabbitMQ routes message to durable queue \texttt{sales.hq}.
  \item HO listener consumes the message and upserts into \texttt{consolidated\_sales}.
\end{enumerate}

For delete, the same pipeline executes with \texttt{operation=DELETE}; HO removes the matching logical key.

\subsection{Reliability Observations}
Current implementation includes:
\begin{itemize}[leftmargin=*]
  \item Durable exchange and queue.
  \item Persistent message publishing.
  \item Consumer-side idempotent merge key.
\end{itemize}

Potential future hardening:
\begin{itemize}[leftmargin=*]
  \item Publisher confirms and retries.
  \item Dead-letter queue for poison messages.
  \item Explicit monitoring/metrics for lag and failure rates.
\end{itemize}

\section{Known Trade-Offs and Future Enhancements}
\subsection{Trade-Offs}
\begin{itemize}[leftmargin=*]
  \item With direct publish, full transactional outbox guarantees are not present.
  \item If broker is unavailable at exact publish instant, application-level recovery strategy is limited unless additional retry/persistence logic is introduced.
\end{itemize}

\subsection{Recommended Enhancements}
\begin{enumerate}[leftmargin=*]
  \item Add publisher confirms and failure callbacks.
  \item Introduce dead-letter and retry policies at queue level.
  \item Add integration tests for failure scenarios and idempotent replay.
  \item If strong guaranteed delivery is required, reintroduce an outbox pattern with event relay worker, but with optimized logging and polling strategy.
\end{enumerate}

\section{Conclusion}
The implemented system provides a clear, operational, and technically coherent distributed synchronization architecture:
\begin{itemize}[leftmargin=*]
  \item Branches remain independently operable and expose local CRUD UIs/APIs.
  \item RabbitMQ mediates asynchronous BO-to-HO propagation.
  \item HO applies deterministic upsert/delete logic using a stable composite identity.
  \item The direct-publish redesign simplifies runtime behavior and removes scheduler polling noise while preserving event-driven consolidation semantics.
\end{itemize}

In summary, RabbitMQ is not only a transport utility in this implementation; it is the core integration mechanism that allows the system to remain modular, asynchronous, and extensible while maintaining coherent data synchronization across organizational boundaries.

\end{document}
